Las aplicaciones web modernas exponen interfaces de programación cuya superficie de ataque incluye la inyección SQL, el \textit{Cross-Site Scripting}, el \textit{Path Traversal} y la inyección de comandos del sistema operativo. Estas categorías figuran entre los riesgos más críticos del estándar OWASP Top~10 \parencite{owasp2021} y comparten una propiedad: la evidencia del ataque reside en una sola petición HTTP. Las startups y empresas tecnológicas que desarrollan servicios sobre Node.js rara vez cuentan con un equipo de seguridad dedicado. Este trabajo parte de la premisa de que, para esas organizaciones, las soluciones perimetrales resultan costosas de operar y mantener. De ahí surge la necesidad de una herramienta de protección que opere dentro de la propia aplicación.

Este trabajo presenta logSguarDian, una librería npm que implementa un middleware de autoprotección de aplicaciones en tiempo de ejecución (RASP) para Express. El middleware intercepta cada petición, extrae un vector de 76 características numéricas y lo evalúa con un modelo híbrido. Un Random Forest clasifica la petición en cinco clases (tráfico benigno y las cuatro categorías de ataque) y decide el bloqueo, y un Isolation Forest asigna un puntaje de anomalía que se registra y no bloquea. Ambos modelos se serializan en formato ONNX y se ejecutan en \texttt{worker\_threads}, para no bloquear el \textit{Event Loop} del proceso anfitrión. El proyecto se desarrolló junto con Sebastián Huertas, quien colaboró en el componente de ciberseguridad y en la implementación del lado npm de la librería, y este documento reporta el componente de aprendizaje automático.

El objetivo general es diseñar e implementar esa librería para que detecte y clasifique las amenazas cubiertas, y marque anomalías estadísticas, sin comprometer el rendimiento del servidor anfitrión. Tres objetivos específicos lo desarrollan. El primero fija un límite de 5~ms por petición para el pipeline de extracción de características. El segundo establece umbrales de F1, de \textit{recall} y de tasa de falsos positivos para el modelo híbrido sobre un conjunto de prueba bloqueado. El tercero exige una cobertura mínima de 80\,\% por categoría de ataque y un incremento del p95 de latencia de hasta 15~ms en un entorno de prueba controlado. Este trabajo no formula una hipótesis estadística; los criterios de aceptación de los tres objetivos específicos, definidos en el capítulo de Objetivos, cumplen esa función.

La investigación es aplicada y experimental. El corpus reúne 383\,202 filas de 17 fuentes, de las cuales 268\,240 se destinan al entrenamiento: nueve conjuntos o listas públicas, siete generadores propios con tráfico sintético, creados para cubrir vacíos de cobertura, y una fuente de telemetría curada. La evaluación comprende métricas del modelo sobre un conjunto de prueba bloqueado, cuatro configuraciones de integración sobre una aplicación vulnerable (sin protección, solo logSguarDian, solo un cortafuegos de aplicaciones web y ambos en capas) y pruebas en cuatro aplicaciones de terceros. Las cifras principales se acompañan de su procedencia y de las limitaciones que las afectan.

Los resultados cumplen, con los matices que detalla el capítulo de Discusión, los criterios del extractor, del modelo híbrido y de la cobertura por categoría, aunque con el par A (tarball e09120\ldots) el intervalo de confianza de la cobertura de \texttt{path\_traversal} incluye valores inferiores al 80\,\%; cumplen también el de latencia en carga normal y de ataque. No cumplen los criterios de latencia en carga de volumen, de memoria adicional ni de CPU. El trabajo documenta además las limitaciones que acotan la lectura de las métricas: la lectura repetida del conjunto de prueba, la estructura de casi-duplicados del corpus y el peso del corpus benigno en los falsos positivos observados en aplicaciones de terceros.

El documento se organiza como sigue. Los capítulos de Antecedentes, Justificación, Objetivos y Alcance plantean el problema y delimitan el trabajo, y el Marco Teórico desarrolla los conceptos necesarios. La Metodología describe las fases, las técnicas y los criterios de evaluación. El capítulo de Resultados presenta las mediciones, y el de Discusión las interpreta. Los capítulos de Conclusiones y Recomendaciones cierran el trabajo, y los anexos reúnen los análisis complementarios y la procedencia de los artefactos.
